\section{R10 执行记录：显式 fragment 映射的真实 K2 写回}

\subsection{实验目标与边界}

R9 已经在最小 kernel 中证明了 K2 的 K\_INTER/\code{SM90\_U32x4\_STSM\_N}
映射不能直接套用到 row-major global tile，同时给出了逐元素 exact 的逆映射。
R10 将该公式接入真实 \FlashKDA K2 recurrence。实验只在 B300/SM103a 上进行，
所有新路径均为编译期 opt-in，默认构建仍然使用原有 output shared-memory/TMA
epilogue。

本轮分为三个阶段：先验证 scalar direct global store 的完整正确性，再把相邻
BF16 元素打包为 32-bit store，最后和同一源码的 V-split control 及默认四个
MMA warp 路径比较 CUDA-event 延迟与 NCU resource 指标。实验成功的判据是
output、final state 在所有覆盖形状上逐 bit exact；性能只有在正确性通过后才
用于决策。

\subsection{生产 kernel 的 opt-in 修改}

\paragraph{编译开关。}

\code{FlashKDA/setup.py} 新增两个环境变量到 NVCC 定义：

\begin{itemize}
  \item \code{FLASH\_KDA\_C1\_DIRECT\_OUTPUT} 定义
        \code{C1\_K2\_DIRECT\_OUTPUT=1}，启用 scalar 显式写回；
  \item \code{FLASH\_KDA\_C1\_DIRECT\_OUTPUT\_VEC} 定义
        \code{C1\_K2\_DIRECT\_OUTPUT\_VEC=1}，启用相邻 BF16 pair 的
        32-bit 打包写回。
\end{itemize}

两个宏默认都是零，因此未设置环境变量的构建不会改变已有二进制语义。
\code{C1\_K2\_DIRECT\_OUTPUT\_VEC} 只有在 direct output 已经打开时才有
实际作用。

\paragraph{scalar direct path。}

在 MMA warp 完成 Phase 4 的 \code{out\_bf16[i]} 后，不再把 fragment 写入
shared output tile，而是按 R9 公式计算每个元素的全局坐标。设
\(j\) 是经 \code{smem\_thr\_store\_C.retile\_S} 后的 fragment index，
则：

\begin{lstlisting}
row       = (lane / 4) + (((j / 2) & 1) * 8)
local_col = ((j / 4) * 8) + ((lane & 3) * 2) + (j & 1)
tile_col  = (warp * 2 + block) * 16 + local_col
col       = value_slice * (D / 2) + tile_col  // ValueSplit
col       = tile_col                         // non-split
\end{lstlisting}

对每个实际 token 使用 \code{row < actual\_len} 谓词；全局线性地址为
\code{(bos + t * CHUNK + row) * H * D + head\_idx * D + col}。ValueSplit
时两个 CTA 分别负责不相交的 64 列，因此不会发生 CTA 间写冲突。

direct path 同时跳过 output producer/consumer pipeline 和 STORE warp 的
output TMA transaction。state update 仍由 MMA warp 执行；在最终 state store
前增加 CTA-wide \code{\_\_syncthreads()}，替代原来由 output pipeline wait
提供的发布顺序。partial tail 不写入下一 tile 的邻居区域。

\paragraph{pair-packed path。}

一个 fragment 的 8 个元素恰好形成四个连续 BF16 pair：\(j=\{0,1\}\)、
\(\{2,3\}\)、\(\{4,5\}\)、\(\{6,7\}\)。R10-C 读取两个元素的 raw
BF16 storage，组成

\begin{lstlisting}
packed = uint32_t(src(c0).storage) |
         (uint32_t(src(c1).storage) << 16);
*reinterpret_cast<uint32_t*>(out_raw_ptr + global_base + col) = packed;
\end{lstlisting}

每个 pair 的起始列是偶数，且 K2 的 head/value-slice 基址按 4-byte 对齐，
所以该版本不改变 BF16 bit pattern，只把四次 16-bit store 降为四次 32-bit
store。scalar 版本保留在同一开关下，作为地址和边界谓词的对照。

\subsection{构建过程与一次无效实验}

远端工作副本为
\code{/home/lcpu/10578328/wmhpc-training-camp-x-lcpu-ai-infra-seminars}。
第一次 job 24533 使用 UV 隔离构建，失败原因是临时 build environment 找不到
PyTorch：

\begin{lstlisting}
ModuleNotFoundError: No module named 'torch'
\end{lstlisting}

后续统一使用 \code{uv pip install --reinstall --no-build-isolation -e}，并在
GPU allocation 内设置 \code{FLASH\_KDA\_CUDA\_ARCHS=103a}。随后发现第一次
同步把头文件放到了 \code{FlashKDA/fwd\_kernel2.cuh}，而 setuptools 实际
编译的是 \code{FlashKDA/csrc/smxx/fwd\_kernel2.cuh}。因此 job 24534--24543
期间的 direct/control 构建实际使用旧源码；对应的 1.16 ms benchmark 只说明
两个 binary 相同，不能用于判断 R10 优化，日志仍保留在
\code{r10\_direct\_build\_24540.log}、\code{r10\_direct\_bench\_fixed\_h96\_b300\_24539.log}
等文件中。

修正同步目标后，实际构建命令为：

\begin{lstlisting}
cd /home/lcpu/10578328/wmhpc-training-camp-x-lcpu-ai-infra-seminars/assignment02
srun --partition=gpu --gres=gpu:1 --cpus-per-task=32 \\
  bash -lc 'FLASH_KDA_C1_VSPLIT_K2=1 \\
    FLASH_KDA_C1_DIRECT_OUTPUT=1 \\
    FLASH_KDA_CUDA_ARCHS=103a \\
    FLASH_KDA_VERSION_SUFFIX=+r10.direct2 \\
    uv pip install --reinstall --no-build-isolation \\
    -e team/c1_flashkda/FlashKDA'
\end{lstlisting}

pair-packed 版本额外设置
\code{FLASH\_KDA\_C1\_DIRECT\_OUTPUT\_VEC=1}；默认对照清除
\code{FLASH\_KDA\_C1\_VSPLIT\_K2}、\code{FLASH\_KDA\_C1\_DIRECT\_OUTPUT}
和 \code{FLASH\_KDA\_C1\_DIRECT\_OUTPUT\_VEC}。成功构建日志包括
\code{r10\_direct\_build\_24545.log}、\code{r10\_vec\_build\_24551.log}、
\code{r10\_full\_directvec\_build\_24560.log} 和
\code{r10\_ncu\_base\_build\_24568.log}。

\subsection{R10-A：scalar direct output 正确性}

基础 probe \code{vsplit\_native\_smoke.py} 覆盖 5 个 state I/O 组合；扩展
probe 额外覆盖固定长度跨 chunk、varlen、\(B=2\)、BF16/FP32 state。V-split
和 non-split 两种编译形态都得到相同结果：

\begin{table}[htbp]
\centering
\small
\begin{tabular}{l r r r r}
\toprule
variant & basic cases & extended cases & output mismatch & state mismatch \\
\midrule
V-split scalar & 5 & 5 & 0 & 0 \\
non-split scalar/vector build & 5 & 5 & 0 & 0 \\
\bottomrule
\end{tabular}
\caption{R10 direct path 的 exactness 汇总。每个 case 均同时检查 output 和适用的 final state。}
\end{table}

代表性结果来自 job 24546 的 scalar smoke、job 24585 的 scalar extended smoke
以及 job 24557/24563 的 pair-packed 扩展 smoke：\(T=16,H=96\)、\(T=97,H=96\)、varlen
\([17,33,65]\)、\(B=2,T=17,H=4\)、\(B=2,T=64,H=4\) 均为
\code{output\_different=0}、\code{state\_different=0}。这说明显式映射、
tail 谓词、ValueSplit 列偏移和 state publication barrier 在完整 K2 中均能
保持 reference 的 BF16 bit pattern；没有只针对 R9 诊断值成立的假 exactness。

\subsection{R10-C：B300 性能对照}

所有正式 event benchmark 使用 warmup=20、iterations=100、repeats=5，
输入为固定 \([8192,96,128]\)。下表只比较同一源码、同一节点上确认过编译宏
的结果；job 24539--24544 的错误同步结果不纳入。

\begin{table}[htbp]
\centering
\small
\begin{tabular}{l r r r}
\toprule
variant & BF16 state (ms) & no state (ms) & FP32 state (ms) \\
\midrule
V-split control & 2.0291 & 2.0280 & 2.0505 \\
V-split scalar direct & 2.1088 & 2.0659 & 2.0762 \\
V-split pair-packed direct & 1.1744 & 1.1591 & 1.1553 \\
non-split default & 1.0285 & 1.0285 & 0.9979 \\
non-split pair-packed direct & 1.6779 & 1.7037 & 1.6529 \\
\bottomrule
\end{tabular}
\caption{B300 R10 fixed H96 CUDA-event benchmark。}
\end{table}

scalar direct 相对 V-split control 慢约 \(3.9\%\)，说明显式整数地址计算和
逐元素 store 的成本超过了省下的 output pipeline。pair-packed 将 V-split
延迟降到 \(1.1744\) ms，相对 V-split control 减少约 \(42.1\%\)，证明
“减少输出写回阶段”在两个 CTA 的 split 路径上确实有效。但 V-split control
本身不是当前最佳路径；pair-packed direct 仍比非 split default 的
\(1.0285\) ms 慢约 \(14.2\%\)。将同一 map 用于四个 MMA warp 的 non-split
kernel 后反而为 \(1.6779\) ms，比 default 慢约 \(63.1\%\)。

varlen 的 non-split pair-packed 结果也没有反转结论：job 24565 在
\code{[1300,547,2048,963,271,3063]} 和八段 1024 的输入上分别得到
BF16-state \(1.4910\) ms 和 \(1.2140\) ms。它们均 exact，但没有超过
现有 default/官方路径的相同形状记录。

\subsection{资源分析：direct store 没有减少寄存器或 shared memory}

对 non-split default 和 pair-packed direct 分别运行 NCU basic set。NCU
replay 时间不作为性能数据，只读取 Launch Statistics、Occupancy 和
DRAM throughput：

\begin{table}[htbp]
\centering
\small
\begin{tabular}{l r r r r}
\toprule
variant & block & registers/thread & dynamic smem/block (B) & achieved occupancy \\
\midrule
default & 192 & 66--74 & 98,432 & 9.36--9.38\% \\
pair-packed direct & 192 & 78--92 & 98,432 & 9.37--9.38\% \\
\bottomrule
\end{tabular}
\caption{NCU basic set 的 recurrence resource 结果；范围来自不同 state I/O 实例化。}
\end{table}

direct 版本的寄存器数不是下降，而是增加约 5--26 registers/thread；NCU 的
Dynamic Shared Memory Per Block 在两者均为 98,432 B，因为本轮没有删除 union
中预留的 output storage。早期记录的 200,704 B 是 Shared Memory Configuration/
SM，不可误作单个 CTA 的 dynamic shared-memory 数。
两个版本的 occupancy 几乎相同。direct capture 的 DRAM throughput 约为
11.3--11.9\%，default 约为 10.8--11.4\%，没有显示出“少一次 TMA 就显著
降低内存压力”的证据，反而说明普通 global store 和地址计算增加了活动。
对应原始 CSV 为
\code{r10\_full\_directvec\_ncu\_basic\_24567.csv} 和
\code{r10\_full\_base\_ncu\_basic\_24569.csv}。

\subsection{R10 判定}

本轮完成了 R10-A 的生产集成和 R10-C 的 pair-packed 尝试，但没有产生可替换
当前默认路径的端到端优化：

\begin{itemize}
  \item 正确性：scalar 和 pair-packed 在 V-split/non-split、固定长度、
        varlen、tail、batch、BF16/FP32 state 上均 exact；
  \item split 局部收益：pair-packed direct 比 V-split control 快约
        \(42.1\%\)，证明 R9 显式 map 不是纯诊断产物；
  \item 端到端目标：当前最佳 non-split default 为 \(1.0285\) ms，
        pair-packed non-split 为 \(1.6779\) ms，未超过 baseline；
  \item 资源解释：direct path 增加寄存器、保留相同 shared-memory 配置，
        因而没有形成预期的 occupancy/resource 优势。
\end{itemize}

因此 \code{C1\_K2\_DIRECT\_OUTPUT} 和
\code{C1\_K2\_DIRECT\_OUTPUT\_VEC} 均保持 opt-in，不进入默认构建。
R10 的 direct-store 主线到此停止继续堆叠普通 global-store 变体。下一轮应
转向删除 output storage 的实际 shared-memory 占用、重新设计 warp-specialized
epilogue，或分析 launch/TMA 固定开销；若不能同时保持当前 non-split 的
低寄存器路径，继续增加显式地址计算只会扩大代码和寄存器压力。

\begin{record}{R10 结论}
R10 证明了 R9 显式输出映射可以安全接入真实 K2：两个编译形态和完整
exactness 矩阵均通过。pair-packed 映射显著改善了 V-split control，但
non-split default 仍快 \(14.2\%\)，且 NCU 显示 direct path 寄存器更多、
shared memory 不变。故本轮的有效成果是“可复用的 exact direct epilogue
和清晰的负性能结论”，不是默认 kernel 加速；默认环境保持不变。
\end{record}
